

%%%%%%%%%%%%%%%%%%%%%%%%%%%%%%%%%%%%%%%%%%%%%%%%%%%%%%%%%%%%%
\section{模型的评价}

\subsection{模型的优点}

\begin{itemize}
  \item \textbf{优点一：方法论系统完整。}本文按照从数据探索到归因分析、再到预测建模和业务应用的四阶段递进框架，构建了完整的零售销售数据分析与预测体系，每个阶段的方法选择有明确的理论依据和实际考量。
  \item \textbf{优点二：统计分析与业务逻辑深度融合。}问题二中引入ANOVA效应量（$\eta^2$）量化各因素的影响力，既保证了统计推断的严谨性，又通过方差解释比例实现了业务的直观解读，识别出产品维度是销售额差异的第一驱动力。
  \item \textbf{优点三：模型选择务实高效。}问题三在多模型对比中并未盲目追求复杂算法，而是根据小样本（24个月训练数据）的实际情况，选择了Holt-Winters指数平滑这一经典但适用性极强的方法，以较低的模型复杂度获得了最优的预测效果（MAPE=23.56\%，$R^2$=0.72）。
  \item \textbf{优点四：策略建议具有数据支撑和可操作性。}问题四的五个维度经营建议均基于前文的定量分析结果（预测数据、效应量排序、增长趋势），每条建议均有具体的数值依据，避免了泛泛而谈。
\end{itemize}

\subsection{模型的缺点}

\begin{itemize}
  \item \textbf{缺点一：预测精度受限于训练数据量。}MAPE为23.56\%，虽然在小样本条件下可以接受，但对于库存管理和供应链规划等对精度要求较高的场景仍有提升空间。主要瓶颈在于训练数据仅有4年（48个月），限制了模型对长期趋势和异常年份的捕捉能力。
  \item \textbf{缺点二：模型未充分利用外部协变量。}Holt-Winters模型仅依赖历史销售额本身，未融入宏观经济指标、节假日效应、促销活动等外部信息，导致模型在面对结构性突变或异常事件时的预测能力有限，2025年Q3-Q4的预测偏差即源于此。
\end{itemize}
